\begin{definition}
  The local contribution of a prime exponent is $w(0)=1$, $w(1)=-1$, and $w(k)=0$ for $k\ge2$.

  \begin{lean}
    def mobiusWeight (k : ℕ) : ℤ := if k = 0 then 1 else if k = 1 then -1 else 0
  \end{lean}
\end{definition}

\begin{definition}
  Define the M\"obius function for $n>0$ by multiplying $w(e_n(p))$ over the prime support.
  Thus $\mu(1)=1$, a prime square forces value zero, and a product of $r$ distinct primes has value
  $(-1)^r$.
  Set $\mu(0)=0$ in the total definition.

  \begin{lean}
    noncomputable def mobius (n : ℕ) : ℤ :=
      if n = 0 then 0 else (primeExponents n).prod (fun _ k => mobiusWeight k)
  \end{lean}
\end{definition}

\begin{theorem}
  \label{ent:thm:mobius-one}
  The M\"obius function has $\mu(1)=1$.

  \begin{lean}
    theorem mobius_one : mobius 1 = 1
  \end{lean}
\end{theorem}

\begin{proof}
  The prime support of one is empty, so the defining product is one.

  \begin{lean}
    simp only [mobius, one_ne_zero, ↓reduceIte, primeExponents_one, Finsupp.prod_zero_index]
  \end{lean}
\end{proof}

\begin{theorem}
  \label{ent:thm:mobius-prime-pow}
  For a prime $p$, $\mu(p^k)=w(k)$: the values are one for $k=0$, minus one for $k=1$, and zero
  thereafter.

  \begin{lean}
    theorem mobius_prime_pow (p k : ℕ) (hp : Chapter02.isPrime p) :
        mobius (p ^ k) = mobiusWeight k
  \end{lean}
\end{theorem}

\begin{proof}
  The exponent vector of $p^k$ has the single coordinate $k$ at $p$.
  Evaluate the defining product.

  \begin{lean}
    have hpos : 0 < p ^ k := pow_pos (by
        have := hp.1
        omega) k
    rw [mobius, ite_eq_right (ne_of_gt hpos), primeExponents_prime_pow p k hp]
    exact Finsupp.prod_single_index (by simp only [mobiusWeight, ↓reduceIte])
  \end{lean}
\end{proof}

\begin{theorem}
  \label{ent:thm:mobius-eq-zero-of-prime-sq-dvd}
  If $n>0$ and $p^2\mid n$ for a prime $p$, then $\mu(n)=0$.

  \begin{lean}
    theorem mobius_eq_zero_of_prime_sq_dvd (n p : ℕ) (hn : 0 < n)
        (hp : Chapter02.isPrime p) (hdiv : p ^ 2 ∣ n) : mobius n = 0
  \end{lean}
\end{theorem}

\begin{proof}
  The divisibility criterion gives $e_n(p)\ge2$.
  Its contribution is zero, so the whole product vanishes.

  \begin{lean}
    have hle := (dvd_iff_primeExponents_le (p ^ 2) n
      (pow_pos (by
          have := hp.1
          omega) 2) hn).1 hdiv
    rw [primeExponents_prime_pow p 2 hp, Finsupp.single_le_iff] at hle
    rw [mobius, ite_eq_right (ne_of_gt hn), Finsupp.prod]

    have hmem : p ∈ (primeExponents n).support :=
      Finsupp.mem_support_iff.mpr (by omega)

    apply Finset.prod_eq_zero hmem
    simp only [mobiusWeight, ite_eq_right (by omega : primeExponents n p ≠ 0),
      ite_eq_right (by omega : primeExponents n p ≠ 1)]
  \end{lean}
\end{proof}

\begin{theorem}
  \label{ent:thm:mobius-eq-sign}
  If positive $n$ has no prime-square divisor and has $r$ distinct prime factors, then
  $\mu(n)=(-1)^r$.

  \begin{lean}
    theorem mobius_eq_sign (n : ℕ) (hn : 0 < n)
        (hsq : ∀ p : ℕ, Chapter02.isPrime p → ¬p ^ 2 ∣ n) :
        mobius n = (-1 : ℤ) ^ (primeExponents n).support.card
  \end{lean}
\end{theorem}

\begin{proof}
  Every exponent in the prime support is positive and less than two, hence equals one.
  Every factor in the defining product is therefore minus one.

  \begin{lean}
    rw [mobius, ite_eq_right (ne_of_gt hn), Finsupp.prod]

    have hweight : ∀ p ∈ (primeExponents n).support, mobiusWeight (primeExponents n p) = -1 := by
      intro p hps

      have hp := (primeExponents_spec n hn).1 p hps

      have hpos := Nat.pos_of_ne_zero (Finsupp.mem_support_iff.mp hps)

      have hlt : primeExponents n p < 2 := by
        by_contra hh
        apply hsq p hp
        apply (dvd_iff_primeExponents_le _ n (pow_pos (by
            have := hp.1
            omega) 2) hn).2
        rw [primeExponents_prime_pow p 2 hp, Finsupp.single_le_iff]
        omega

      have he : primeExponents n p = 1 := by omega
      simp only [he, mobiusWeight, one_ne_zero, ↓reduceIte]
    rw [Finset.prod_congr rfl hweight, Finset.prod_const]
  \end{lean}
\end{proof}

\begin{theorem}
  \label{ent:thm:mobius-multiplicative}
  The M\"obius function is multiplicative.

  \begin{lean}
    theorem mobius_multiplicative : IsMultiplicative mobius
  \end{lean}
\end{theorem}

\begin{proof}
  Coprime integers have disjoint prime supports: a common prime would divide their gcd, one.
  Multiplication adds exponent vectors, so the defining product splits into the two disjoint
  products.

  \begin{lean}
    intro m n hm hn hcop
    rw [mobius, mobius, mobius, ite_eq_right (ne_of_gt (mul_pos hm hn)),
      ite_eq_right (ne_of_gt hm), ite_eq_right (ne_of_gt hn), primeExponents_mul m n hm hn]

    exact Finsupp.prod_add_index_of_disjoint (prime_support_disjoint m n hm hn hcop) _
  \end{lean}
\end{proof}

\begin{theorem}
  \label{ent:thm:sum-mobius}
  For positive $n$, \[ \sum_{d\mid n}\mu(d)=
    \begin{cases}
      1 & n=1, \\0&n>1.
    \end{cases} \]

  \begin{lean}
    theorem sum_mobius (n : ℕ) (hn : 0 < n) :
        divisorSum mobius n = if n = 1 then 1 else 0
  \end{lean}
\end{theorem}

\begin{proof}
  The divisor sum of $\mu$ is multiplicative.
  On every positive prime power it is $1-1+0+\cdots+0=0$.
  For $n>1$, its nonempty prime-power factorization therefore has a zero factor.
  At one the sole summand is one.

  \begin{lean}
    by_cases he : n = 1

    · subst n
      simp only [divisorSum, Nat.divisors_one, Finset.sum_singleton, mobius_one, ↓reduceIte]

    · rw [ite_eq_right he, multiplicative_prime_powers _
        (divisorSum_multiplicative _ mobius_multiplicative) n (by omega)]

      have hs : (primeExponents n).support.Nonempty := by
        by_contra h

        have hh := (primeExponents_spec n hn).2
        change (∏ p ∈ (primeExponents n).support, p ^ primeExponents n p) = n at hh
        rw [Finset.not_nonempty_iff_eq_empty.mp h, Finset.prod_empty] at hh
        exact he hh.symm

      obtain ⟨p, hp⟩ := hs
      apply Finset.prod_eq_zero hp
      exact divisorSum_mobius_prime_pow p _ ((primeExponents_spec n hn).1 p hp)
        (Nat.pos_of_ne_zero (Finsupp.mem_support_iff.mp hp))
  \end{lean}
\end{proof}

\begin{theorem}
  \label{ent:thm:mobius-inversion}
  Let $F(n)=\sum_{d\mid n}f(d)$ for every positive $n$.
  Then \[ f(n)=\sum_{d\mid n}\mu(d)F(n/d). \]

  \begin{lean}
    theorem mobius_inversion (f F : ℕ → ℤ)
        (hF : ∀ n : ℕ, 0 < n → F n = divisorSum f n) (n : ℕ) (hn : 0 < n) :
        f n = ∑ d ∈ n.divisors, mobius d * F (n / d)
  \end{lean}
\end{theorem}

\begin{proof}
  Extend the functions by zero at zero.
  Dirichlet convolution is the finite sum over pairs with product $n$, and multiplication by the
  constant-one arithmetic function is divisor summation.

  \begin{lean}
    have hF' : arithmeticFunction F = arithmeticFunction (divisorSum f) := by
      ext k
      by_cases hk : k = 0

      · subst k
        rfl

      · rw [arithmeticFunction_apply _ k (Nat.pos_of_ne_zero hk),
          arithmeticFunction_apply _ k (Nat.pos_of_ne_zero hk), hF k (Nat.pos_of_ne_zero hk)]
  \end{lean}

  The proved cancellation identity says that convolution of $\mu$ with the constant-one function
  is the convolution identity.
  Associativity is just regrouping the finite sum over triples with product $n$.
  Therefore $\mu*F=\mu*(f*1)=f$.
  Expanding at $n$ gives the stated formula.

  \begin{lean}
    have hconv : arithmeticFunction mobius * arithmeticFunction F = arithmeticFunction f := by
      rw [hF', arithmeticFunction_divisorSum, mul_left_comm, mobius_mul_zeta, mul_one]

    have h := congrArg (fun g : ArithmeticFunction ℤ => g n) hconv
    rw [convolution_apply _ _ n hn, arithmeticFunction_apply f n hn] at h

    exact h.symm
  \end{lean}
\end{proof}

\begin{theorem}
  \label{ent:thm:mobius-inversion-complement}
  Under the same hypotheses, \[ f(n)=\sum_{d\mid n}\mu(n/d)F(d). \]

  \begin{lean}
    theorem mobius_inversion' (f F : ℕ → ℤ)
        (hF : ∀ n : ℕ, 0 < n → F n = divisorSum f n) (n : ℕ) (hn : 0 < n) :
        f n = ∑ d ∈ n.divisors, mobius (n / d) * F d
  \end{lean}
\end{theorem}

\begin{proof}
  Interchange the two factors in the convolution, or equivalently permute the divisors by
  $d\mapsto n/d$.

  \begin{lean}
    rw [mobius_inversion f F hF n hn, ← convolution_apply _ _ n hn, mul_comm,
      convolution_apply _ _ n hn]

    apply Finset.sum_congr rfl
    intro d hd
    exact mul_comm _ _
  \end{lean}
\end{proof}

\begin{theorem}
  \label{ent:thm:multiplicative-of-divisorsum}
  If the divisor sum of $f$ is multiplicative, then $f$ is multiplicative.

  \begin{lean}
    theorem multiplicative_of_divisorSum (f : ℕ → ℤ)
        (hf : IsMultiplicative (divisorSum f)) : IsMultiplicative f
  \end{lean}
\end{theorem}

\begin{proof}
  M\"obius inversion writes $f$ as the convolution of two multiplicative functions.
  To see that this convolution is multiplicative, split divisors of a coprime product as $ab$ and
  their complementary divisors as $(m/a)(n/b)$; the double sum factors.
  The checked auxiliary lemma carries out this reindexing.

  \begin{lean}
    have hconv := convolution_multiplicative mobius (divisorSum f) mobius_multiplicative hf

    intro m n hm hn hcop

    have hh := hconv m n hm hn hcop
    dsimp only at hh

    rw [mobius_inversion f (divisorSum f) (fun _ _ => rfl) (m * n) (mul_pos hm hn),
      hh,
      ← mobius_inversion f (divisorSum f) (fun _ _ => rfl) m hm,
      ← mobius_inversion f (divisorSum f) (fun _ _ => rfl) n hn]
  \end{lean}
\end{proof}

\begin{theorem}
  \label{ent:thm:multiplicative-iff-divisorsum}
  An arithmetic function is multiplicative if and only if its divisor sum is multiplicative.

  \begin{lean}
    theorem multiplicative_iff_divisorSum (f : ℕ → ℤ) :
        IsMultiplicative f ↔ IsMultiplicative (divisorSum f)
  \end{lean}
\end{theorem}

\begin{proof}
  Combine preservation of multiplicativity with the converse just proved.

  \begin{lean}
    ⟨divisorSum_multiplicative f, multiplicative_of_divisorSum f⟩
  \end{lean}
\end{proof}

\begin{definition}
  The summatory M\"obius function is $M(N)=\sum_{k=1}^N\mu(k)$.
  It is the number of square-free integers in this interval with an even number of prime factors,
  minus the number with an odd number; one belongs to the even group.

  \begin{lean}
    noncomputable def mertens (N : ℕ) : ℤ := ∑ k ∈ Finset.Icc 1 N, mobius k
  \end{lean}
\end{definition}
